%% Whitepaper-only short form of the two-source appendix.
%% The full version, with the calibration argument written out, is
%% appendix_twosource_body.tex, which sa-ai.tex \inputs.
%% Figures from repro/analytic/regen_twosource_appendix.py (analytic only).

The rate uses the single even coordinate $\hat\Omega\hat I$, which represents
both the inviscid inflectional and the viscous Tollmien--Schlichting mechanisms
because the two co-vary across the Falkner--Skan family. They separate on
the parabola, where $\hat I\equiv0$: plane Poiseuille returns zero rate at
every Reynolds number against a true $Re_c=5772$ (Orszag 1971), and on the
favorable branch the model's amplification is an order of magnitude below
the Drela--Giles envelope by $\beta\!=\!0.35$ (Fig.~\ref{fig:tscalibrate}).
Blasius has no inflection point, so the missing branch is the main mechanism,
not an exotic one.

A second coordinate separates them:
\begin{equation}
  P_{\mathrm{c}} \;=\; \hat\Omega\,\big\langle -\hat Z\big\rangle_+,
  \qquad \hat Z = Z/U_\mathrm{loc},
  \label{eq:ts_pcurv}
\end{equation}
positive wherever $u''<0$, clipped at an adverse wall, vanishing like $d$ at
the wall and zero in the free stream. Each branch has its own rate and its own onset, and the two \emph{sources} are combined:
\begin{gather}
  P_{\mathrm{AI}} = \omega\,\tilde\nu\;
  \operatorname{softmax}_2\!\Big[\,
     a_{\max}\,P_I\;\sigma_{Re}\big(Re_\Omega/Re_\Omega^{\mathrm c,i}\big),\;\;
     a_{\mathrm{visc}}\,P_{\mathrm{c}}\;\sigma_{Re}\big(Re_\Omega/Re_\Omega^{\mathrm c,v}\big)
  \Big],
  \label{eq:ts_twosource}\\[3pt]
  Re_\Omega^{\mathrm c,i} = k\operatorname{softmin}_2\!\Big(C,\,A+B\,P_I^{-2}\Big),
  \qquad
  Re_\Omega^{\mathrm c,v} = A + B_{\mathrm c}\,P_{\mathrm{c}}^{-2},
  \qquad
  \operatorname{softmax}_2(x,y)=\sqrt{x^2+y^2}.
  \notag
\end{gather}
The inviscid gate is unchanged, ceiling included; the viscous branch shares
its floor $A$ and adds one coefficient $B_{\mathrm c}$. Constants:
$a_{\max}\!=\!0.19$, $a_{\mathrm{visc}}\!=\!0.0276$, $A\!=\!124.6$,
$B\!=\!1.424$, $B_{\mathrm c}\!=\!130$, $C\!=\!1851.2$.

The gates are separate, not soft-minimized into one. On Blasius the two
thresholds are $358$ and $2598$; sharing the lower switches the viscous rate
on early and changes the Blasius anchor by $+7.4\%$. With them separate the
construction reproduces the single-source rate to within $0.7\%$ over the
calibrated family (Table~\ref{tab:tsfamily}), so $a_{\max}$, $k$ and the
airfoil campaign are untouched. Dropping the ceiling $C$ reduces $\beta\!=\!+0.20$ to $0.63\times$ its single-source value.

$a_{\mathrm{visc}}$ is fitted to the Drela--Giles envelope on the favorable
Falkner--Skan family. Plane Poiseuille then fixes $B_{\mathrm c}$
independently and without refitting, since $P_I\equiv0$ on a parabola:
requiring the frozen-profile eigenvalue of Eq.~\eqref{eq:ts_twosource} to
change sign at $Re_c=5772$ gives Table~\ref{tab:tschannel}. At
$a_{\mathrm{visc}}=0.0276$ the channel requires $B_{\mathrm c}=91$ against the
family's $130$, and the family's constants place the channel neutral point
at $7409$, $1.28\times$ the true value. The channel is a calibration and
verification device only; its observed transition is subcritical and
finite-amplitude.

Three limits. $P_{\mathrm c}$ is largest at a stagnation line, where the
filled attachment profile keeps $\langle-\hat Z\rangle_+$ finite while
$P_I\!\le\!0$ keeps the single-source rate at zero at the nose, so
the branch would amplify in the one region where the present formulation does
not; adopting it requires either showing its
stagnation-line rate stays below onset at leading-edge Reynolds numbers or
pairing it with a nose treatment. The viscous rate
$a_{\mathrm{visc}}P_{\mathrm{c}}$ is Reynolds-independent, inherited from
the inviscid branch. And the branches decompose the rate, not the physics:
the viscous gate on Blasius opens only near $Re_\theta\!\approx\!1190$, so
the Blasius rate remains that of the $P_I$ branch. The branch is not part of the model
validated here.

\begin{figure}[tp]\centering
\includegraphics[width=0.99\textwidth]{twosource_calibrate.pdf}
\caption{Two-source rate against Drela--Giles across the Falkner--Skan
family, $\beta=+1.0$ to $-0.1988$; analytic march, no CFD.
(a)~rate as early ($N\!\in\![1,5]$) and late ($N\!\in\![5,9]$) secants, with
the single-source rate. (b)~marched $N\!=\!1$ station against the
Drela--Giles $N\!=\!1$ station $Re_{\theta0}+1/(dN/dRe_\theta)$; the diagonal
is exact agreement.}\label{fig:tscalibrate}
\end{figure}

\begin{figure}[tp]\centering
\includegraphics[width=0.99\textwidth]{twosource_onsetgraze.pdf}
\caption{Each onset branch against its own family, at each profile's
Drela--Giles critical $Re_{\theta0}(H)$. (a)~inflectional branch
$A+B/P_I^2$, adverse family. (b)~curvature branch
$A+B_{\mathrm c}/P_{\mathrm c}^2$, low-$H$ members, in the coordinate
$P_{\mathrm c}$ of Eq.~\eqref{eq:ts_pcurv}.}\label{fig:tsgraze}
\end{figure}

\begin{figure}[tp]\centering
\includegraphics[width=0.72\textwidth]{twosource_transport.pdf}
\caption{Disturbance transport under Eq.~\eqref{eq:ts_twosource} on four
wedges---adverse, Blasius, favorable, and stagnation
($\beta\!=\!+1$)---against each profile's Drela--Giles envelope. On the
stagnation wedge the branch switches on at
$Re_\theta\!\approx\!2.5\times10^3$ against $Re_{\theta0}\!=\!6.5\times10^3$.
Same instrument and same $c_{\nu,\mathrm{ai}}=1/6$ as the model's own
calibration.}\label{fig:tstransport}
\end{figure}

\begin{figure}[tp]\centering
\includegraphics[width=0.99\textwidth]{twosource_sphere.pdf}
\caption{Low-$H$ Falkner--Skan family on the $\mathbb{RP}^2$ indicator
sphere, $H\!\approx\!2.22$ to Blasius $H\!=\!2.59$, with (b) near-wall
detail. The family bunches against the neutral parabola locus where $\hat I$
is too small to separate members whose Drela--Giles critical Reynolds
numbers differ by nearly a factor of two; $P_{\mathrm c}$ separates
them.}\label{fig:tssphere}
\end{figure}

\begin{table}[tp]
  \centering\small
  \caption{Peak dimensionless source $a\,\sigma_{Re}$ on Falkner--Skan
  wedges: single-source, and Eq.~\eqref{eq:ts_twosource} at
  $(a_{\mathrm{visc}},B_{\mathrm c})=(0.0276,130)$.}
  \label{tab:tsfamily}
  \begin{tabular}{rrccc}
    \toprule
    $\beta$ & $Re_\theta$ & single-source & two-source & ratio\\
    \midrule
    $0$        & $500$  & $0.01485$ & $0.01485$ & $1.000$\\
    $0$        & $1000$ & $0.01485$ & $0.01487$ & $1.001$\\
    $-0.10$    & $400$  & $0.03082$ & $0.03083$ & $1.000$\\
    $-0.1988$  & $300$  & $0.09692$ & $0.09764$ & $1.007$\\
    $+0.10$    & $1000$ & $0.00695$ & $0.00696$ & $1.000$\\
    $+0.20$    & $2000$ & $0.00281$ & $0.00282$ & $1.004$\\
    $+0.35$    & $4000$ & $0.00033$ & $0.00405$ & $12.4$\\
    \bottomrule
  \end{tabular}
\end{table}

\begin{table}[tp]
  \centering\small
  \caption{Plane-channel check: the $B_{\mathrm c}$ placing the neutral point
  on $Re_c=5772$, and the neutral point obtained at the
  Falkner--Skan-anchored $B_{\mathrm c}=130$.}
  \label{tab:tschannel}
  \begin{tabular}{ccc}
    \toprule
    $a_{\mathrm{visc}}$ & $B_{\mathrm c}$ for $Re_c\!=\!5772$ & $Re_c$ at $B_{\mathrm c}\!=\!130$\\
    \midrule
    $0.0200$ & $81.0$  & $8095$\\
    $0.0276$ & $91.0$  & $7409$\\
    $0.0350$ & $99.9$  & $6916$\\
    $0.0600$ & $127.9$ & $5834$\\
    \bottomrule
  \end{tabular}
\end{table}
